\section{Geospatial format and engine benchmarks}
\label{related-work:geospatial-format-and-engine-benchmarks}

An early portable benchmark for spatial database engines is the Jackpine suite of \textcite{ray2011_jackpine_spatial_database_benchmark}, designed to fill the gap left by the absence of a TPC-Spatial standard. Jackpine combines micro queries based on the \acrshort{ogc} DE-9IM topological model with six macro workload scenarios, including map search and browsing, geocoding, flood-risk analysis, land information management, and toxic-spill analysis. The benchmark was applied to PostgreSQL, MySQL, and Informix on TIGER data for Texas, combined with City of Austin and Travis County GIS datasets, with the test harness running on the same machine as the database server. The benchmark targets single-node engines and predates the current cloud-native formats.

For distributed spatial systems, \textcite{pandey2018_how_good_modern_spatial_analytics} compared five Spark-based systems on Amazon \acrshort{emr} clusters of \num{1} to \num{16} nodes. The systems evaluated were SpatialSpark, GeoSpark, Magellan, Simba, and LocationSpark, run against \acrfull{osm} data comprising approximately \num{500} million records across point, line-string, rectangle, and polygon datasets. The study evaluated range, $k$-nearest-neighbor, distance-join, spatial-join, and $k$NN-join queries, reporting query throughput, total run-time, peak execution memory, and intra-cluster shuffle read and write volumes. GeoSpark, later renamed Apache Sedona, supported the broadest set of data types and queries and delivered the best performance in most cases, although it lacked $k$NN-join support. End-to-end network transfer between client and cluster was not measured.

Empirical comparisons of cloud-native vector formats against traditional formats appear mainly in community and vendor reports. \textcite{holmes2023_geoparquet_performance} compared GeoParquet, Shapefile, GeoPackage, and FlatGeobuf on the Google Open Buildings dataset, measuring write times and file sizes for inputs ranging from \num{498} megabytes to \num{21} gigabytes. GeoParquet produced the smallest files in every case, while Shapefile failed on the largest dataset. \textcite{flatgeobuf2024_flatgeobuf_benchmarks} reported relative read times for FlatGeobuf and Shapefile on Danish \acrshort{osm} roads (\num{906602} linestrings) and Danish cadastral polygons (\num{2511772} polygons), with FlatGeobuf reading the full datasets \num{2.2} and \num{8.3} times faster than Shapefile, respectively. Both studies use local disk storage and report only timing and file size; neither measures read or transfer over a network.

Two limitations recur across this literature. First, network transfer is rarely measured. \textcite{tang2020_locationspark_distributed_spatial_query} treats network communication cost as a system-design concern that motivates the spatial bitmap filter optimization rather than as a reported experimental metric, and \textcite{folkerts2012_benchmarking_in_the_cloud} argues that cloud benchmarks must report end-to-end performance jointly with pricing, of which data transfer is a major component. This thesis reports bytes received and bytes sent as primary metrics and derives the corresponding cost, addressing both axes directly. Second, single-node and distributed engines are rarely compared in the same study; \textcite{pandey2018_how_good_modern_spatial_analytics} compares only Spark-based systems, while Jackpine and similar single-node benchmarks predate the Spark-based spatial systems that emerged in the mid-2010s.